\section{A polynomial bound for the complete refinement cost}
\label{sec:refinement_output_polynomial}

The following numerical estimate combines both reference profiles, the
sample size and the final refinement ceiling. Its coefficient and
integer exponent precede length, multiplicity and requested accuracy.

\begin{lemma}[A common monomial for the recorded refinement cost]
    \label{lem:refinement_output_polynomial}
    \lean{Real.exists_refinement_output_polynomial_bound}
    \leanok
    \uses{lem:reference_accuracy_log_profile,lem:reference_profile_polynomial,
        lem:viable_refinement_costs_le_polynomial,lem:viable_refinement_error_parameters}
    Fix $\gamma>0$, $0<g\leq1$, $u,K_{
m tree},K_{
m ref}\geq0$
    and an integer $D\geq1$. There are $C>0$ and $k\in\mathbb N$,
    depending only on these fixed data, with the following property.
    Let $L\geq0$, $r\geq1$ be integers, $\eta>0$ and
    $\gamma\leq A\leq L+1$. Put
    \begin{align}
        c&=\gamma/A,&
        \rho&=\frac{\eta\gamma}{2A+\eta\gamma},&
        \varepsilon&=c^2\rho/32768,\\
        b_{
m tree}&=16+\log(1/g)+\log\bigl(320000r(L+1)\bigr),\\
        b_{
m ref}&=16+\log(1/g)+\log\bigl(32r(L+1)/\varepsilon\bigr),\\
        F(b)&=\bigl(u^3/g+g^{-1/4}b^{3/4}u\bigr)(1+\log b).
    \end{align}
    Let integers $B_{
m tree},B_{
m ref}\geq1$ satisfy
    $\log B_{
m tree}\leq K_{
m tree}F(b_{
m tree})$ and
    $\log B_{
m ref}\leq K_{
m ref}F(b_{
m ref})$.
    If the integer $s$ obeys
    $s\leq4097r(1+\log r+\log D+\log2)$, define
    \begin{align}
        B_*&=\max\{r^2B_{
m tree},r^2B_{
m ref}\},&
        N&=\left\lceil\frac{100D^{12}}{c\rho}\right\rceil,&
        M&=D^6s^3.
    \end{align}
    Then, with $Z=(L+1)(r+1)(1+\eta^{-1})$,
    \begin{align}
        10^6(B_*N)^2(L+1)M^2\leq CZ^k.
    \end{align}
\end{lemma}
\begin{proof}\leanok
    Put $J=1+\gamma^{-1}$, $K_c=6J^2$ and
    $K_\varepsilon=196608J^3$. Since $A,\eta^{-1}\leq Z$ and $Z\geq1$,
    Lemma~\ref{lem:viable_refinement_costs_le_polynomial} gives
    $(c\rho)^{-1}\leq K_cZ^3$ and
    $\varepsilon^{-1}\leq K_\varepsilon Z^4$.
    Also $r(L+1)\leq Z$, so both accuracy logarithms lie between one
    and $b_0+\log(Z^5)$, where
    $b_0=16+\log(1/g)+\log(32K_\varepsilon)$.

    Set $K=\max\{K_{
m tree},K_{
m ref}\}$,
    $a_0=u^3/g$, $a_1=g^{-1/4}u$,
    $p_0=\lceil K(a_0+5a_1)\rceil$ and
    $C_b=e^{K(a_0+5a_1)b_0}$.
    Monotonicity of the profiles and
    Lemma~\ref{lem:reference_bound_le_nat_power}, applied at $X=Z^5$,
    give $B_{
m tree},B_{
m ref}\leq C_bZ^{5p_0}$.
    Thus $B_*\leq C_bZ^{5p_0+2}$.

    For $C_s=4097(1+\log D+\log2)$,
    the inequality $\log r\leq r-1$ gives $s\leq C_sr^2$.
    Consequently $M\leq C_MZ^6$, where $C_M=D^6C_s^3$.
    The ceiling estimate gives $N\leq C_NZ^3$, where
    $C_N=100D^{12}K_c+1$.
    Hence the recorded cost is at most
    $10^6(C_bC_N)^2C_M^2Z^{10p_0+23}$.
    One may take $k=10p_0+23$ and
    $C=1+10^6(C_bC_N)^2C_M^2$.
\end{proof}

This is a numerical consequence of the reference bounds used in
\cite[lines 1777--1817]{Arad2016RigorousRG}. It does not itself construct
a physical subspace or establish the computational cost of a procedure.


\begin{lemma}[The logarithmic profile of a polynomial accuracy budget]
    \label{lem:reference_accuracy_log_profile}
    \lean{Real.reference_accuracy_log_profile_bounds}
    \leanok
    Let $0<g\le1$, $K,Z,r,l\ge1$, $rl\le Z$ and
    $0<\varepsilon\le1$. Suppose $\varepsilon^{-1}\le KZ^4$ and put
    \begin{align}
        b&=16+\log(1/g)+\log(32rl/\varepsilon),&
        b_0&=16+\log(1/g)+\log(32K).
        \label{eq:reference_accuracy_profile}
    \end{align}
    Then $1\le b\le b_0+\log(Z^5)$.
\end{lemma}
\begin{proof}\leanok
    The logarithm arguments are at least one. Moreover
    $32rl/\varepsilon\le32KZ^5$; take logarithms and split the product.
\end{proof}

\begin{lemma}[A monotone reference profile gives a natural-power bound]
    \label{lem:reference_profile_polynomial}
    \lean{Real.reference_profile_polynomial_bound}
    \leanok
    Let $0\le K_0\le K$, $a_0,a_1\ge0$, $b,b_0,X\ge1$ and $B>0$.
    Suppose $b\le b_0+\log X$ and
    $\log B\le K_0(a_0+a_1b^{3/4})(1+\log b)$. Then
    \begin{align}
        B\le e^{K(a_0+5a_1)b_0}
          X^{\lceil K(a_0+5a_1)\rceil}.
        \label{eq:reference_profile_polynomial}
    \end{align}
\end{lemma}
\begin{proof}\leanok
    \uses{lem:reference_bound_le_nat_power}
    The powers and logarithms are monotone on the positive half-line.
    Replace $b$ by $b_0+\log X$ and $K_0$ by $K$, then apply the
    natural-power reference bound.
\end{proof}
